% ENEM 2020-2023 — 11 questões MCQ — Aritmética I
% Fonte: lista "Aritmética I" do site projetoagathaedu.com.br
% Omitidas: Q1,Q4,Q6,Q7,Q12 (imagens), Q17 (conteúdo truncado)

---
tipo: Múltipla Escolha
dificuldade: Média
nivel: médio
disciplina: Matemática
assunto: Sistemas de Equações
gabarito: B
tags: [ENEM, sistema de equações, metrô, tiquete, aritmética]
fonte: concurso
concurso: ENEM
banca: INEP
ano: 2023
numero: 2
---
\question
O metrô de um município oferece dois tipos de tiquetes com colorações diferentes, azul e vermelha, sendo vendidos em cartelas, cada qual com nove tiquetes da mesma cor e mesmo valor unitário. Duas cartelas de tiquetes azuis e uma cartela de tiquetes vermelhos são vendidas por R\$ 32,40. Sabe-se que o preço de um tiquete azul menos o preço de um tiquete vermelho é igual ao preço de um tiquete vermelho mais cinco centavos.

Qual o preço, em real, de uma cartela de tiquetes vermelhos?
\begin{choices}
\choice 4,68
\CorrectChoice 6,30
\choice 9,30
\choice 10,50
\choice 10,65
\end{choices}

---
tipo: Múltipla Escolha
dificuldade: Média
nivel: médio
disciplina: Matemática
assunto: Grandezas e Medidas
gabarito: A
tags: [ENEM, grandezas, unidades, radiação, sievert, conversão]
fonte: concurso
concurso: ENEM
banca: INEP
ano: 2023
numero: 3
---
\question
Os raios cósmicos são fontes de radiação ionizante potencialmente perigosas para o organismo humano. Para quantificar a dose de radiação recebida, utiliza-se o sievert (Sv), definido como a unidade de energia recebida por unidade de massa. A exposição à radiação proveniente de raios cósmicos aumenta com a altitude, o que pode representar um problema para as tripulações de aeronaves. Recentemente, foram realizadas medições acuradas das doses de radiação ionizante para voos entre Rio de Janeiro e Roma. Os resultados têm indicado que a dose média de radiação recebida na fase de cruzeiro (que geralmente representa 80\% do tempo total de voo) desse trecho intercontinental é $2\ \mu$Sv/h. As normas internacionais da aviação civil limitam em 1\,000 horas por ano o tempo de trabalho para as tripulações que atuem em voos intercontinentais. Considere que a dose de radiação ionizante para uma radiografia torácica é estimada em 0,2 mSv.

A quantas radiografias torácicas corresponde a dose de radiação ionizante à qual um tripulante que atue no trecho Rio de Janeiro--Roma é exposto ao longo de um ano?
\begin{choices}
\CorrectChoice 8
\choice 10
\choice 80
\choice 100
\choice 1000
\end{choices}

---
tipo: Múltipla Escolha
dificuldade: Média
nivel: médio
disciplina: Matemática
assunto: Frações
gabarito: C
tags: [ENEM, frações, proporção, sociedade, capital]
fonte: concurso
concurso: ENEM
banca: INEP
ano: 2020
numero: 5
---
\question
Antônio, Joaquim e José são sócios de uma empresa cujo capital é dividido, entre os três, em partes proporcionais a 4, 6 e 6, respectivamente. Com a intenção de igualar a participação dos três sócios no capital da empresa, Antônio pretende adquirir uma fração do capital de cada um dos outros dois sócios.

A fração do capital de cada sócio que Antônio deverá adquirir é
\begin{choices}
\choice $\dfrac{1}{2}$
\choice $\dfrac{1}{3}$
\CorrectChoice $\dfrac{1}{9}$
\choice $\dfrac{2}{3}$
\choice $\dfrac{4}{3}$
\end{choices}

---
tipo: Múltipla Escolha
dificuldade: Fácil
nivel: médio
disciplina: Matemática
assunto: Sistemas de Equações
gabarito: B
tags: [ENEM, sistema de equações, frações, infrações, trânsito]
fonte: concurso
concurso: ENEM
banca: INEP
ano: 2020
numero: 8
---
\question
Em um país, as infrações de trânsito são classificadas de acordo com sua gravidade. Infrações dos tipos leves e médias acrescentam, respectivamente, 3 e 4 pontos na carteira de habilitação do infrator, além de multas a serem pagas. Um motorista cometeu 5 infrações de trânsito. Em consequência, teve 17 pontos acrescentados em sua carteira de habilitação.

Qual é a razão entre o número de infrações do tipo leve e o número de infrações do tipo média cometidas por esse motorista?
\begin{choices}
\choice $\dfrac{1}{4}$
\CorrectChoice $\dfrac{3}{2}$
\choice $\dfrac{3}{4}$
\choice $\dfrac{5}{17}$
\choice $\dfrac{7}{17}$
\end{choices}

---
tipo: Múltipla Escolha
dificuldade: Fácil
nivel: médio
disciplina: Matemática
assunto: Grandezas e Medidas
gabarito: A
tags: [ENEM, grandezas, tempo, conversão, minutos, segundos]
fonte: concurso
concurso: ENEM
banca: INEP
ano: 2020
numero: 9
---
\question
Os tempos gastos por três alunos para resolver um mesmo exercício de matemática foram: 3,25 minutos; 3,4 minutos e 191 segundos.

O tempo gasto a mais, em segundos, pelo aluno que concluiu por último a resolução do exercício, em relação ao primeiro que o finalizou, foi igual a
\begin{choices}
\CorrectChoice 13.
\choice 14.
\choice 15.
\choice 21.
\choice 29.
\end{choices}

---
tipo: Múltipla Escolha
dificuldade: Fácil
nivel: médio
disciplina: Matemática
assunto: Grandezas e Medidas
gabarito: D
tags: [ENEM, grandezas, unidades, volume, comprimento, conversão]
fonte: concurso
concurso: ENEM
banca: INEP
ano: 2020
numero: 10
---
\question
Três pessoas, X, Y e Z, compraram plantas ornamentais de uma mesma espécie que serão cultivadas em vasos de diferentes tamanhos.

O vaso escolhido pela pessoa X tem capacidade de 4 dm³. O vaso da pessoa Y tem capacidade de 7\,000 cm³ e o de Z tem capacidade igual a 20 L.

Após um tempo do plantio das mudas, um botânico que acompanha o desenvolvimento delas realizou algumas medições e registrou que a planta que está no vaso da pessoa X tem 0,6 m de altura. Já as plantas que estão nos vasos de Y e Z têm, respectivamente, alturas medindo 120 cm e 900 mm.

O vaso de maior capacidade e a planta de maior altura são, respectivamente, os de
\begin{choices}
\choice Y e X.
\choice Y e Z.
\choice Z e X.
\CorrectChoice Z e Y.
\choice Z e Z.
\end{choices}

---
tipo: Múltipla Escolha
dificuldade: Fácil
nivel: médio
disciplina: Matemática
assunto: Grandezas e Medidas
gabarito: E
tags: [ENEM, grandezas, volume, conversão, litro, mililitro]
fonte: concurso
concurso: ENEM
banca: INEP
ano: 2020
numero: 11
---
\question
É comum as cooperativas venderem seus produtos a diversos estabelecimentos. Uma cooperativa láctea destinou 4 m³ de leite, do total produzido, para análise em um laboratório da região, separados igualmente em 4\,000 embalagens de mesma capacidade.

Qual o volume de leite, em mililitro, contido em cada embalagem?
\begin{choices}
\choice 0,1
\choice 1,0
\choice 10,0
\choice 100,0
\CorrectChoice 1\,000,0
\end{choices}

---
tipo: Múltipla Escolha
dificuldade: Média
nivel: médio
disciplina: Matemática
assunto: Grandezas e Medidas
gabarito: A
tags: [ENEM, grandezas, velocidade média, tempo, distância]
fonte: concurso
concurso: ENEM
banca: INEP
ano: 2020
numero: 13
---
\question
Um motorista fez uma viagem de 100 km partindo da cidade A até a cidade B. Nos primeiros 30 km, a velocidade média na qual esse motorista viajou foi de 90 km/h. No segundo trecho, de 40 km, a velocidade média foi de 80 km/h. Suponha que a viagem foi realizada em 1 h 30 min.

A velocidade média do motorista, em quilômetro por hora, no último trecho da viagem foi de
\begin{choices}
\CorrectChoice 45.
\choice 67.
\choice 77.
\choice 85.
\choice 113.
\end{choices}

---
tipo: Múltipla Escolha
dificuldade: Fácil
nivel: médio
disciplina: Matemática
assunto: Grandezas e Medidas
gabarito: E
tags: [ENEM, grandezas, volume, tempo, litro, consumo]
fonte: concurso
concurso: ENEM
banca: INEP
ano: 2020
numero: 14
---
\question
Considere que cada uma das cinco pessoas de uma família toma dois banhos por dia, de 15 minutos cada. Sabe-se que a cada hora de banho são gastos aproximadamente 540 litros de água. Considerando que um mês tem 30 dias,

a quantidade total de litros de água consumida, nos banhos dessa família, durante um mês, é mais próxima de
\begin{choices}
\choice 1\,350.
\choice 2\,700.
\choice 20\,250.
\choice 20\,520.
\CorrectChoice 40\,500.
\end{choices}

---
tipo: Múltipla Escolha
dificuldade: Média
nivel: médio
disciplina: Matemática
assunto: Frações
gabarito: E
tags: [ENEM, frações, aritmética, custo, amendoim, desconto]
fonte: concurso
concurso: ENEM
banca: INEP
ano: 2020
numero: 15
---
\question
A fim de reforçar o orçamento familiar, uma dona de casa começou a produzir doces para revender. Cada receita é composta de $\frac{4}{5}$ de quilograma de amendoim e $\frac{1}{5}$ de quilograma de açúcar.

O quilograma de amendoim custa R\$ 10,00 e o do açúcar, R\$ 2,00. Porém, o açúcar teve um aumento e o quilograma passou a custar R\$ 2,20. Para manter o mesmo custo com a produção de uma receita, essa dona de casa terá que negociar um desconto com o fornecedor de amendoim.

Nas condições estabelecidas, o novo valor do quilograma de amendoim deverá ser igual a
\begin{choices}
\choice R\$ 9,20.
\choice R\$ 9,75.
\choice R\$ 9,80.
\choice R\$ 9,84.
\CorrectChoice R\$ 9,95.
\end{choices}

---
tipo: Múltipla Escolha
dificuldade: Difícil
nivel: médio
disciplina: Matemática
assunto: Grandezas e Medidas
gabarito: B
tags: [ENEM, grandezas, tempo, futebol, cronômetro, operações]
fonte: concurso
concurso: ENEM
banca: INEP
ano: 2020
numero: 16
---
\question
Uma partida de futebol tem dois tempos de 45 minutos cada. A duração do intervalo entre cada tempo é de 15 minutos. No Brasil, o segundo tempo é iniciado zerando-se o cronômetro, mas em campeonatos europeus, começa com o cronômetro posicionado em 45 minutos. Em uma partida de um campeonato europeu, um time marcou um gol aos 17 minutos e 45 segundos. A outra equipe empatou o jogo aos 54 minutos e 32 segundos. O tempo do intervalo foi respeitado e houve um acréscimo de 2 minutos ao primeiro tempo do jogo.

O tempo transcorrido entre os dois gols foi de
\begin{choices}
\choice 54 minutos e 47 segundos.
\CorrectChoice 53 minutos e 47 segundos.
\choice 51 minutos e 47 segundos.
\choice 38 minutos e 47 segundos.
\choice 36 minutos e 47 segundos.
\end{choices}
